\begin{statementbox}
% Exact ellipse x^2/9+y^2/4=1; chord y=2x/3.
\begin{center}
\begin{tikzpicture}[x=1cm,y=1cm,font=\small]
 \path[use as bounding box] (-3.8,-3.4) rectangle (5.1,3);
 \node at (0,2.7) {椭圆斜弦：整弦与两个坐标投影};
 \draw[->,gray] (-3.5,0)--(3.5,0) node[right] {$x$};
 \draw[->,gray] (0,-2.1)--(0,2.3) node[left] {$y$};
 \draw[blue!60!black,thick] (0,0) ellipse (3 and 2);
 \coordinate (A) at ({-3/sqrt(2)},{-2/sqrt(2)});
 \coordinate (B) at ({3/sqrt(2)},{2/sqrt(2)});
 \coordinate (H) at ({3/sqrt(2)},{-2/sqrt(2)});
 \draw[very thick,magenta!70!black] (A)--(B) node[midway,above,sloped] {$|AB|$};
 \draw[thick,teal!70!black,dashed] (A)--(H)
 --(B);
 \draw[teal!70!black,thin] ({3/sqrt(2)},.65)--(3.25,.65);
 \node[right] at (3.3,.65) {$|y_1-y_2|$};
 \draw (H) ++(-.18,0)--++(0,.18)--++(.18,0);
 \fill (A) circle (1.6pt) node[left] {$A$};
 \fill (B) circle (1.6pt) node[right] {$B$};
 \node[below right] at (H) {$H$};
 \draw[teal!70!black,dotted] (A)--({-3/sqrt(2)},-2.2) (H)--({3/sqrt(2)},-2.2);
 \draw[<->,teal!70!black] ({-3/sqrt(2)},-2.2)--({3/sqrt(2)},-2.2) node[midway,below] {$|x_1-x_2|$};
 \node at (0,-3.05) {$k=\frac23,\quad |AB|=\sqrt{1+k^2}\,|x_1-x_2|$};
\end{tikzpicture}
\end{center}


\textbf{初中知识入口（圆的弦）：}
圆的半径为 $r$、圆心到弦 $AB$ 的距离为 $d$（$0\le d<r$）时，
由垂径定理与勾股定理有
$\displaystyle |AB|=2\sqrt{r^2-d^2}$。
这是\textbf{圆的专用计算法}。
推广到圆锥曲线时，改用“坐标投影差 $+$ 勾股定理”构建通式。

\medskip
\textbf{\textcolor{CStatement}{条件}}
\begin{description}[style=nextline,leftmargin=2.8em,labelsep=0.5em,itemsep=0.4em]
  \item[\textbf{条件一（两交点）}] 直线 $l$ 与圆锥曲线交于两个\textbf{不同的实点}
  $A(x_1,y_1)$、$B(x_2,y_2)$。
  \item[\textbf{条件二（斜率）}] 若 $l:y=kx+b$，则斜率 $k$ 存在；
  允许 $k=0$，但不包含竖直直线。
\end{description}

\textbf{\textcolor{CStatement}{核心结论（横坐标形式）}}
\[
  \boxed{|AB|=\sqrt{1+k^2}\,|x_1-x_2|.}
\]
\textbf{直观描述：}实际弦长 $=$ 水平投影长度 $\times\sqrt{1+k^2}$。

\medskip
\textbf{\textcolor{CStatement}{推广结论（纵坐标形式）}}
当 $k\neq0$ 时，
\[
  \boxed{|AB|=\sqrt{1+\frac{1}{k^2}}\,|y_1-y_2|.}
\]
\textbf{直观描述：}实际弦长 $=$ 竖直投影长度 $\times\sqrt{1+1/k^2}$。

\medskip
\textbf{\textcolor{CStatement}{特殊直线与统一形式}}
\begin{itemize}[leftmargin=*,itemsep=0.3em]
 \item 水平线 $y=b$（$k=0$）：$\boxed{|AB|=|x_1-x_2|}$。
 \item 竖直线 $x=c$（$k$ 不存在）：$\boxed{|AB|=|y_1-y_2|}$。
 \item 更实用的纵坐标写法：只要直线写为 $x=ty+d$（\textbf{包括竖直线} $t=0$），
 \[
  \boxed{|AB|=\sqrt{1+t^2}\,|y_1-y_2|.}
 \]
 \item 全方向统一：设直线参数式为 $(x,y)=(x_0,y_0)+s(u,v)$，$(u,v)\ne(0,0)$；两交点参数为 $s_1,s_2$，则
 \[
  \boxed{|AB|=\sqrt{u^2+v^2}\,|s_1-s_2|.}
 \]
\end{itemize}

\textbf{\textcolor{CStatement}{高考速算（判别式形式）}}
若代入 $y=kx+b$ 后得到
$\alpha x^2+\beta x+\gamma=0$，其中 $\alpha\ne0$、
$\Delta=\beta^2-4\alpha\gamma>0$，则
\[
 \boxed{|AB|=\frac{\sqrt{1+k^2}\sqrt{\Delta}}{|\alpha|}.}
\]
\textbf{注意：}只有确实得到一元\textbf{二次}方程，才能使用此判别式速算式；不得将退化为一次方程的情形误判成两交点。
\end{statementbox}
